\RequirePackage[OT1]{fontenc}
\documentclass[10pt,journal,onecolumn]{IEEEtran}
\usepackage{graphicx,amsmath,amssymb,booktabs,array,url,longtable,needspace}
\setlength{\LTpre}{6pt}
\setlength{\LTpost}{6pt}
\usepackage[nocompress]{cite}
\usepackage[colorlinks=true,allcolors=blue]{hyperref}
\renewcommand{\thesection}{S-\Roman{section}}
\renewcommand{\thesectiondis}{S-\Roman{section}.}
\renewcommand{\thesubsection}{\thesection-\Alph{subsection}}
\renewcommand{\thetable}{S-\Roman{table}}
\renewcommand{\thefigure}{S\arabic{figure}}
\newcommand{\ARC}{ARC-Audit}
\title{ARC-Audit: Protocol, Security, and Experimental Supplement}
\author{Anonymous Author(s)}
\begin{document}
\bstctlcite{IEEEexample:BSTcontrol}
\maketitle
This supplement accompanies \emph{ARC-Audit: Dependency-Aware Auditing
and Authenticated Record Recovery for Industrial IoT Data}. It specifies
the encrypted object format, dependency interfaces, proof premises, and the
complete measurements underlying the three research questions. The evidence
uses two 21,600-record windows for construction and the complete train1/train2
segments for auditing and recovery. Tables retain the corresponding units
and comparison populations.

\section{Encrypted Objects and Public Geometry}\label{sec:protocol}
\subsection{Canonical Representation and Disclosure}
The logical SDA-GGD container stores private schema values, constants,
record time/label content, shared bases, references, and deviations. Public
geometry records the encoded length, record count, label offsets and widths,
and group modes, widths, base-bit positions, dictionary cardinalities and
stream offsets. This geometry permits addressing of complete ciphertext
objects without decrypting BaseIDs. Public hashes commit to ciphertext.

For a logical bit range $[a,a+w)$, the first and last physical objects are
$\lfloor a/(8\cdot4068)\rfloor$ and
$\lfloor(a+w-1)/(8\cdot4068)\rfloor$ when $w>0$. Empty ranges map to an empty
set. The map is applied independently within each batch; an object is named
by batch identity and local index. A page may contain portions of several
logical streams. The complete page is authenticated whenever any required
range intersects it.

The initial exact-recovery set contains private framing and target reference
ranges. Target label bits and local deviation/RAW content are included in
the final dependency set. Each GD BaseID is authenticated and decrypted before
its value is used to select a dictionary row. A RAW group contributes local
content without a dictionary lookup.

\subsection{AEAD Framing and Key Material}
Each construction samples a 256-bit batch identifier and encryption master
secret. HKDF-SHA256 uses the batch as salt and \texttt{ARC/ENC/v1} as context
to derive a 32-byte key~\cite{RFC5869}. Page index supplies the 12-byte nonce.
AES-256-GCM encrypts at most 4068 plaintext bytes, producing a ciphertext
with a 16-byte tag~\cite{RFC5116}. The stored frame prepends the nonce and
has at most 4096 bytes.

Associated data contains protocol, owner domain, batch, object identity,
codec and plaintext length. These fields are serialized canonically. The
manifest descriptor adds complete frame length and SHA-256 digest. For a
logical container of $B_F>0$ bytes, the object count is
$N=\lceil B_F/4068\rceil$, and $B_{\rm cipher}=B_F+28N$.
Virtual zeros used in audit sectorization or RS coding are not stored as
ciphertext padding. The final object's true length is authenticated.

The manifest-signing and public-audit secrets are independently generated.
The signing public key and expected manifest digest are provisioned as 64 binary
bytes of client pin material. Owner-side files also contain the encryption
secret, batch-scoped base catalog, private framing and base locators.
The local-disk metric includes the actual serialized owner files, whereas
local-base bytes count only catalog content.

\subsection{Manifest and Object Identity}
The public manifest records format, batch, domain, plaintext/ciphertext page
capacities, layout, descriptors, public audit parameters, tag digest and
audit readiness. Ciphertext objects are created before their audit tags;
the complete manifest is then canonically serialized and signed. The audit
identity binds protocol, parameter fingerprint and the entire object descriptor.

Manifest verification checks the expected digest, Ed25519 signature,
canonical serialization, format, capacities, expected object count, and
per-object batch/index/length binding. Public parameters are decoded and
validated in their proper groups. Ciphertext verification checks identity,
valid length and committed hash before authenticated decryption.

\section{Algorithm Details}\label{sec:configuration}
\subsection{Deterministic Group Configuration}
Group positions concatenate fields in stored order and bits in increasing
field-bit order. A stable ascending sort of empirical binary entropy
preserves local order on ties. The configurator tests prefixes of this
order as base positions, beginning with the empty basis. It stops after all
positions or when every row has a distinct basis. The first minimum-cost
prefix wins ties; RAW is selected whenever it costs no more than GD.

The cost includes a group header of $11+2|G|+2k$ bytes, a dictionary stream,
a reference stream and a deviation stream, each byte-rounded independently.
The dependency-ratio parameter is the largest basis frequency divided by
row count. The evaluated value $\rho=1$ admits every candidate.

Sample ranking uses $\min(50{,}000,n)$ rows without replacement, with sorted
indices and fixed build seed. Thus each C1 window uses every row for ranking.
The sampled joint gain is computed directly from two separate groups and
their ordered concatenation. Only positive gains enter the max queue.
Equal priorities follow insertion order; stale group versions and rejected
pairs are discarded. A lower refreshed gain may be deferred at most eight
times before full-data examination. A merge requires both strictly positive
serialized gain and strictly positive payload gain on the complete input.
Accepted groups retain their concatenated field order.

\subsection{Verified Opening and Exact Recovery}
The protected fetch helper first checks the requested original ciphertext.
If a protected object is absent or invalid, it gathers distinct valid stripe
slots. Original objects are checked against their parent descriptors; parity
is checked against the signed sidecar digest. Declared virtual zeros occupy
the unused data slots of an incomplete stripe. Eight valid slots permit RS
reconstruction. Each reconstructed real original is truncated to its committed
length, hash-checked, and cached.

The opening helper then repeats the original descriptor check and performs
AEAD authenticated decryption. Verified plaintext enters a request-local
cache. After structural/reference opening, exact planning selects actual
basis rows plus target-local material. A sparse logical adapter supplies the
private parser, and all decoder bit reads are redirected to cache-checked
spans. The parser's record count, label geometry, group modes, widths, offsets
and base positions must agree with public geometry. Any failed required
check produces failure with an empty decoded result map.

Conservative recovery initially acquires the public closure and then follows
the identical exact-planning and decoding helpers. Hence the comparison
changes the demanded availability set while preserving authentication rules.
The RS helper serves both the initial reference stage and subsequent data stage.

\subsection{Ciphertext Protection Allocation}
The per-batch budget determines the number of admitted stripes, while
RS$(10,8)$ fixes the redundancy within each stripe. Parity and its authenticated
sidecar state are charged against the complete batch ciphertext size.
A reference record's exact dependency set contributes once per object to the
FrequencyD count. Dictionary shuffles all objects and then places
 dictionary-containing objects first; Uniform uses the shuffle directly;
FrequencyD stably orders the shuffled objects by decreasing count. A shared
stripe count is reduced until the largest serialized cost among the three
orders fits the per-batch budget.

For $s_R$ stripes, the charged cost is
\[
B_{\rm side}=8192s_R+|M_R|+64+32+32.
\]
The last three terms are signature, signing public key and digest pin.
$M_R$ records parent snapshot, RS dimensions, page capacity, stripe members
and two parity digests per stripe. Admission uses the actual canonical
serialization length. The budget is
$\lfloor0.10\sum_p\ell(p)\rfloor$. Two 4096-byte parity shards are stored
per stripe, and existing ciphertext originals are reused.

\section{Security Interfaces and Composition}
\subsection{Public-View Confidentiality}
Equal-leakage inputs have identical declared geometry, counts, lengths and
public access schedules. Batch identifiers and secrets are sampled
independently of private content. Under secure derivation, batch keys may
be replaced by independent pseudorandom keys. The nonce-respecting
multi-message AEAD experiment then replaces plaintext slices while retaining
the same public associated data and lengths. This is a hybrid over encrypted
objects and batches.

The reduction generates independent authentication secrets and computes
ciphertext descriptors, audit tags, signatures and RS parity from the
challenge ciphertexts. Public aggregate responses can likewise be generated
from those ciphertext sectors for the public queries. Thus the additional
public artifacts are simulatable post-processing of the encrypted view.
The confidentiality statement uses the specified public leakage; authorized
recovery receives decryption capability as a separate interface.

\subsection{Concrete Backend Interface}\label{sec:backend-interface}
The public authenticator uses additive BLS12-381 groups, 133 sector bases
in $\mathbb G_1$, and a public key in $\mathbb G_2$. Every sector is a
31-byte little-endian integer below the scalar order. The complete frame
is injectively encoded with a signed valid length and prescribed zero padding.
Descriptors and sector bases use distinct hash-to-curve domains~\cite{RFC9380}.

The compact challenge includes the expected snapshot, parameter fingerprint,
targets, full page set, policy/session declaration and fresh nonce.
A canonical-challenge digest enters the coefficient domain. SHA-512 expands
the digest, page index and retry counter; the output is reduced modulo the
scalar order and zero is retried. The deployed coefficients are defined by
this deterministic expansion of a fresh bound challenge.

At the abstract interface, selected-set response authenticity bounds the
probability that any adaptive audit accepts a response with
$\boldsymbol\mu\ne\boldsymbol\mu^\star$, where $\boldsymbol\mu^\star$ is the
specified linear combination of the committed ciphertext sectors. The closure
argument uses this interface property together with authentic metadata.
For the pairing backend used here, the main paper's proposition on selected-set
response authenticity derives the bound under its Type-3 computational
assumption and domain-separated random-oracle model for basis and descriptor
hashing, with distinct basis inputs, page-binding descriptors, and the
expected signed snapshot. Compact coefficients define the checked linear
combination. The reduction applies to the coefficients determined by each
bound challenge and does not require them to be independent uniform draws.
Honest correctness follows separately by substitution into the pairing
equation of the Shacham--Waters authentication structure~\cite{ShachamWaters2013}.

Response parsing enforces framing, dimension, canonical point representation,
subgroup membership and scalar ranges. Public bases are nonidentity points.
An aggregate identity point is parsed within the declared domain and still
must satisfy the pairing equation. The verifier retains its own selected
set; a server-supplied alternative cannot redefine the challenge.

\subsection{Response Serialization}\label{sec:wire-format}
The aggregate response occupies $4+2+48+133\cdot32=4310$ B: protocol framing,
sector count, aggregate point and canonical scalar vector. An individual
response adds an 8-byte page identifier, totaling $4318|U|$ B for a fixed set.
The individual control uses the same fresh coefficient vector and cached
verification context as its aggregate counterpart.

Ciphertext page messages use an 8-byte framing/count prefix and a 12-byte
identifier/length header per returned page. Complete application traffic
also includes the canonical request declaration and serialized challenge.
Owner-MinRef includes its first-stage encrypted page message. A request
spanning $b_T$ data batches accumulates $b_T$ distinct proofs and batch
contexts. Initial manifest/parameter state is charged separately.

\subsection{Binding and Returned Authenticity}
Under authentic pins, a metadata substitution requires a signature or hash
binding failure. Direct and repaired ciphertexts must match the original
signed descriptor. Conditioned on binding, authenticated decryption with
correct keys recovers the original slices, reference parsing selects the
original basis rows, and deterministic decoding reconstructs complete records.
An incorrect returned record therefore requires the binding events defined
in the main paper. Byte withholding instead affects whether the sufficient
availability condition is met.

RS recovery requires eight distinct authentic slots for each missing protected
object. This criterion applies to the true dependency set, including the
initial reference objects. It is a deterministic sufficient condition on an
execution's available objects. The fault experiment samples executions to
measure how frequently recovery succeeds under its stated loss rule.

\section{Workloads and Measurement Definitions}\label{sec:key-evidence}
\subsection{Data and Batches}
The HAI 21.03 train1/train2 source identities are retained in the accompanying
configuration manifests~\cite{Shin2021HAI}. C1 windows each contain 21,600
records; C2 and C3 use 216,001 and 226,801 records. Table~\ref{tab:batches}
gives the batch partition. Each batch is encoded independently and stores
its own ciphertext and authenticated context.

\begin{longtable}{@{}llrr@{}}
\caption{Full-Workload Data Batches}\label{tab:batches}\\
\toprule Segment & Date (2020) & Source offset & Records\\\midrule\endfirsthead
\toprule Segment & Date (2020) & Source offset & Records\\\midrule\endhead
train1 & July 11 & 0 & 86,400\\
train1 & July 12 & 86,400 & 86,400\\
train1 & July 13 & 172,800 & 43,201\\
train2 & July 31 & 0 & 7,200\\
train2 & August 1 & 7,200 & 86,400\\
train2 & August 2 & 93,600 & 86,400\\
train2 & August 3 & 180,000 & 46,801\\\bottomrule
\end{longtable}

The source schema has time, 79 process fields, and four attack indicators.
The independent encoded label stream contains attack\_P1/P2/P3; their
combined indicator is derived. Constants and exact decimal scales belong
to each canonical schema. The C1 windows begin at July 11 00:00 and August 1
00:00, ending at 05:59:59; their source offsets are 0 and 7200.

\subsection{Execution and Timing}
Python is 3.10.12; package versions are NumPy 2.2.6, cryptography 50.0.1,
py-ecc 8.0.0 and the source renderer's Matplotlib 3.10.9. C1 and full-workload
measurement processes use CPU affinity sets $\{2\}$ and $\{4\}$ respectively
with shared host memory and disk. The source manifests bind input digests,
configuration and implementation-file digests.

Block-CZ-256 uses independent Zstd frames at level 3, each covering up to
256 consecutive records. Within a block, nonconstant fields are concatenated
in the fixed schema order, with values retained in source row order and
stored in little-endian byte widths rounded up from their declared bit widths.
The encoder uses a single-threaded compression context with content-size and
checksum fields enabled. The container stores schema and constant metadata,
a separate label stream, and a block index. Record lookup decompresses the
corresponding block frame and combines its field values with this metadata.
The resulting container receives the same page encryption, audit tags, and
signed manifest as the other C1 representations.

C1 builds use seeds 41001 and 41002. Build time is codec configuration and
encoding plus complete secure sealing. Seal time includes encryption, audit
parameter generation, tags, manifest creation and associated file work.
Normalized input preparation and independent full scoring are separate.
Cloud storage counts ciphertext, tags, manifest and its signature; public
parameters reside in the manifest. Base catalog and total owner disk bytes
are reported separately.

C2 uses targets of sizes 1, 16 and 256, with three cases per size/segment.
The seed function binds \texttt{secure-audit}, segment, target count and case.
A request is partitioned among batch contexts and accumulates their measured
local times and bytes. A context is initialized once per batch; values
reused in request rows represent that same initialization. Tables of setup
state therefore describe the contexts needed by requests rather than new
independent initialization repetitions.

C3 uses layout seeds 41001--41005, with 50 cases at each nonzero loss rate.
The seed function binds \texttt{secure-fault}, segment, layout, loss rate,
case and clean status. The first layout supplies 30 clean cases per control.
A signed-sidecar context is initialized before timing the recovery operation.
That initialization is reported separately. The measured recovery operation
includes protected fetching, hash checks, decryption and decoding. External
source comparison follows the service timer.

\subsection{Fault Population and Statistical Units}
Each segment's fault universe contains all original ciphertext objects across
its data batches and the deployed parity slots. The shared admitted stripe
count lets the protected methods use the same slot-loss mask. A target's
retriever sees losses affecting its own batch. NoEC inherits original-object
losses from the same mask. Losses are sampled without replacement with count
$\lceil\alpha N_{\mathrm{fault}}\rceil$, where $N_{\mathrm{fault}}$ is the size
of this fault universe. For train1, $N_{\mathrm{fault}}=5000$ and the three
counts are 5/50/250; for train2, $N_{\mathrm{fault}}=4765$ and they are 5/48/239. The realized fractions
are 0.1000/1.0000/5.0000\% and approximately 0.1049/1.0073/5.0157\%.

Recovery targets exclude FrequencyD reference records: 1440 for the 7200-row
batch and 4096 for each other batch. All eight controls share a case's target
and faults, and execution order rotates across cases. This gives 250 outcomes
per fault point, 12,000 fault outcomes and 480 clean outcomes. Counts remain
in the success denominator whether the operation succeeds or fails.

C1 and C2 report medians with observed minima and maxima ($n=2$ and $3$).
C3 success intervals are pointwise 95\% Wilson intervals. For $k$ successes
in $n$ cases, $\widehat p=k/n$, and $z=1.95996398454$, their center and half-width are
\[
\frac{\widehat p+z^2/(2n)}{1+z^2/n},\qquad
\frac{z\sqrt{\widehat p(1-\widehat p)/n+z^2/(4n^2)}}{1+z^2/n}.
\]
These describe the tested layouts. Success and failure timing summaries use
their respective outcome populations; a missing median denotes an empty
population. Ratios of reported medians are descriptive effect summaries.

\section{C1 Complete Storage and Construction Results}

\begingroup\small\setlength{\tabcolsep}{5pt}
\begin{longtable}{@{}llrrrrrr@{}}
\caption{C1 Secure Storage Ledger (Bytes; Two Builds per Method and Window)}\label{tab:key-c1-statistics}\\
\toprule Segment & Method & Ciphertext & Audit tags & Manifest & Cloud total & Local bases & Local total\\\midrule\endfirsthead
\toprule Segment & Method & Ciphertext & Audit tags & Manifest & Cloud total & Local bases & Local total\\\midrule\endhead
\midrule\multicolumn{8}{r}{Continued on next page}\\\endfoot
\bottomrule\endlastfoot
train1 & SDA-GGD & 1,712,327 & 20,112 & 170,699 & 1,903,138 & 303,595 & 335,884\\
train1 & Per-field-GD & 1,978,108 & 23,184 & 197,520 & 2,198,812 & 208,267 & 245,064\\
train1 & RAW & 13,822,496 & 162,000 & 1,180,606 & 15,165,102 & 0 & 232\\
train1 & Block-CZ-256 & 2,061,529 & 24,192 & 192,280 & 2,278,001 & 0 & 240\\
train2 & SDA-GGD & 1,592,192 & 18,672 & 159,254 & 1,770,118 & 184,687 & 215,702\\
train2 & Per-field-GD & 1,742,643 & 20,448 & 175,748 & 1,938,839 & 174,552 & 209,287\\
train2 & RAW & 13,822,486 & 162,000 & 1,180,606 & 15,165,092 & 0 & 232\\
train2 & Block-CZ-256 & 2,015,313 & 23,664 & 188,382 & 2,227,359 & 0 & 240\\
\end{longtable}
\endgroup


\Needspace{14\baselineskip}
\begingroup\small\setlength{\tabcolsep}{5pt}
\begin{longtable}{@{}llrrrrrr@{}}
\caption{C1 Construction Times (Seconds; Observed Range for Total Build)}\label{tab:c1-time}\\
\toprule Segment & Method & Total & Codec & Tags & Encryption & Min total & Max total\\\midrule\endfirsthead
\toprule Segment & Method & Total & Codec & Tags & Encryption & Min total & Max total\\\midrule\endhead
\midrule\multicolumn{8}{r}{Continued on next page}\\\endfoot
\bottomrule\endlastfoot
train1 & SDA-GGD & 274.356 & 73.661 & 200.419 & 0.00483 & 273.529 & 275.183\\
train1 & Per-field-GD & 231.521 & 1.001 & 230.247 & 0.00537 & 230.801 & 232.240\\
train1 & RAW & 1,620.533 & 0.010 & 1,620.218 & 0.03492 & 1,618.852 & 1,622.214\\
train1 & Block-CZ-256 & 241.711 & 0.030 & 241.401 & 0.00499 & 241.598 & 241.824\\
train2 & SDA-GGD & 235.621 & 47.634 & 187.716 & 0.00427 & 234.945 & 236.296\\
train2 & Per-field-GD & 206.163 & 0.833 & 205.047 & 0.00432 & 206.107 & 206.220\\
train2 & RAW & 1,625.509 & 0.008 & 1,625.187 & 0.03698 & 1,622.530 & 1,628.488\\
train2 & Block-CZ-256 & 236.733 & 0.027 & 236.420 & 0.00472 & 236.020 & 237.445\\
\end{longtable}
\endgroup


Ciphertext includes AEAD framing. Manifest bytes include its signature and public parameters. Local total is the actual owner-directory byte count; local base content is a component of that total. Build time includes the full seal operation, including work beyond the listed component timers.

\section{C2 Complete Audit Measurements}

\begingroup\small\setlength{\tabcolsep}{5pt}
\begin{longtable}{@{}lrlrrrrr@{}}
\caption{C2 Application Bytes and Local Request Time (Seconds)}\label{tab:key-c2-statistics}\\
\toprule Segment & $r$ & Interface & Bytes & Median & Min & Max & $n$\\\midrule\endfirsthead
\toprule Segment & $r$ & Interface & Bytes & Median & Min & Max & $n$\\\midrule\endhead
\midrule\multicolumn{8}{r}{Continued on next page}\\\endfoot
\bottomrule\endlastfoot
train1 & 1 & Public-Aggregate & 6,237 & 3.105910 & 3.076409 & 3.390617 & 3\\
train1 & 1 & Owner-MinRef & 169,601 & 1.163898 & 1.157156 & 1.170814 & 3\\
train1 & 1 & Public-Download & 1,398,654 & 0.003908 & 0.003547 & 0.003921 & 3\\
train1 & 16 & Public-Aggregate & 20,512 & 12.015107 & 11.559540 & 12.017712 & 3\\
train1 & 16 & Owner-MinRef & 2,179,425 & 8.403527 & 8.243958 & 8.552098 & 3\\
train1 & 16 & Public-Download & 5,750,594 & 0.013263 & 0.012822 & 0.013525 & 3\\
train1 & 256 & Public-Aggregate & 35,405 & 32.532051 & 32.345640 & 32.745046 & 3\\
train1 & 256 & Owner-MinRef & 11,163,473 & 31.532080 & 31.203913 & 31.643883 & 3\\
train1 & 256 & Public-Download & 17,041,940 & 0.040020 & 0.039684 & 0.042159 & 3\\
train2 & 1 & Public-Aggregate & 6,254 & 3.169168 & 1.024668 & 3.198473 & 3\\
train2 & 1 & Owner-MinRef & 173,696 & 1.126168 & 0.998525 & 1.178952 & 3\\
train2 & 1 & Public-Download & 1,419,211 & 0.003799 & 0.001079 & 0.003924 & 3\\
train2 & 16 & Public-Aggregate & 19,742 & 10.902284 & 10.556312 & 11.486241 & 3\\
train2 & 16 & Owner-MinRef & 2,138,835 & 8.066283 & 7.674282 & 8.625472 & 3\\
train2 & 16 & Public-Download & 5,018,593 & 0.012236 & 0.011723 & 0.012649 & 3\\
train2 & 256 & Public-Aggregate & 38,524 & 30.908198 & 30.296574 & 31.417652 & 3\\
train2 & 256 & Owner-MinRef & 11,010,612 & 30.148310 & 29.459690 & 30.561948 & 3\\
train2 & 256 & Public-Download & 15,754,868 & 0.035401 & 0.035130 & 0.037838 & 3\\
\end{longtable}
\endgroup


\begingroup\small\setlength{\tabcolsep}{5pt}
\begin{longtable}{@{}lrlrrrrrr@{}}
\caption{C2 Request Accounting and Required Context State (Bytes Except Setup Seconds)}\label{tab:c2-accounting}\\
\toprule Segment & $r$ & Interface & Batches & Response & Stage 1 & Server read & Setup B & Setup s\\\midrule\endfirsthead
\toprule Segment & $r$ & Interface & Batches & Response & Stage 1 & Server read & Setup B & Setup s\\\midrule\endhead
\midrule\multicolumn{9}{r}{Continued on next page}\\\endfoot
\bottomrule\endlastfoot
train1 & 1 & Public-Aggregate & 1 & 4,310 & 0 & 1,392,640 & 666,854 & 8.101\\
train1 & 1 & Owner-MinRef & 1 & 4,310 & 164,328 & 323,584 & 666,854 & 8.101\\
train1 & 1 & Public-Download & 1 & 1,396,728 & 0 & 1,392,640 & 666,854 & 8.101\\
train1 & 16 & Public-Aggregate & 3 & 12,930 & 0 & 5,726,208 & 1,678,681 & 20.816\\
train1 & 16 & Owner-MinRef & 3 & 12,930 & 2,160,832 & 3,633,152 & 1,678,681 & 20.816\\
train1 & 16 & Public-Download & 3 & 5,743,008 & 0 & 5,726,208 & 1,678,681 & 20.816\\
train1 & 256 & Public-Aggregate & 3 & 12,930 & 0 & 16,969,728 & 1,678,681 & 20.816\\
train1 & 256 & Owner-MinRef & 3 & 12,930 & 11,128,596 & 16,408,576 & 1,678,681 & 20.816\\
train1 & 256 & Public-Download & 3 & 17,019,468 & 0 & 16,969,728 & 1,678,681 & 20.816\\
train2 & 1 & Public-Aggregate & 1 & 4,310 & 0 & 1,413,120 & 642,967 & 7.805\\
train2 & 1 & Owner-MinRef & 1 & 4,310 & 168,436 & 319,488 & 642,967 & 7.805\\
train2 & 1 & Public-Download & 1 & 1,417,268 & 0 & 1,413,120 & 642,967 & 7.805\\
train2 & 16 & Public-Aggregate & 3 & 12,930 & 0 & 4,997,120 & 1,554,455 & 19.605\\
train2 & 16 & Owner-MinRef & 3 & 12,930 & 2,115,652 & 3,407,872 & 1,554,455 & 19.605\\
train2 & 16 & Public-Download & 3 & 5,011,784 & 0 & 4,997,120 & 1,554,455 & 19.605\\
train2 & 256 & Public-Aggregate & 4 & 17,240 & 0 & 15,687,596 & 1,623,188 & 21.168\\
train2 & 256 & Owner-MinRef & 4 & 17,240 & 10,972,500 & 15,237,036 & 1,623,188 & 21.168\\
train2 & 256 & Public-Download & 4 & 15,733,588 & 0 & 15,687,596 & 1,623,188 & 21.168\\
\end{longtable}
\endgroup


Values summarize the contexts used by the three request cases. Setup times reused across rows refer to the same actual batch initialization; they are not independent new setup trials. For Public-Download, the response column is the framed ciphertext page message rather than a pairing proof.

\Needspace{18\baselineskip}
\subsection{Representative Mechanism Controls}
The two controls each use one 16-target set, on train1 July 11 and train2 July 31. The page set and fresh coefficients are fixed for the aggregate/individual comparison. The following values are direct control observations.

\begingroup\small\setlength{\tabcolsep}{5pt}
\begin{longtable}{@{}llrr@{}}
\caption{Read Reuse Controls}\label{tab:reuse}\\
\toprule Segment & Traversal & Reads & Bytes\\\midrule\endfirsthead
\toprule Segment & Traversal & Reads & Bytes\\\midrule\endhead
\midrule\multicolumn{4}{r}{Continued on next page}\\\endfoot
\bottomrule\endlastfoot
train1 & Union & 805 & 3,297,280\\
train1 & Separate & 5643 & 23,113,728\\
train1 & Separate-cache & 805 & 3,297,280\\
train2 & Union & 125 & 509,589\\
train2 & Separate & 892 & 3,651,221\\
train2 & Separate-cache & 125 & 509,589\\
\end{longtable}
\endgroup


\begingroup\small\setlength{\tabcolsep}{5pt}
\begin{longtable}{@{}llrrr@{}}
\caption{Aggregation Controls (Bytes and Seconds)}\label{tab:aggregation}\\
\toprule Segment & Response & Proof bytes & Prove s & Verify s\\\midrule\endfirsthead
\toprule Segment & Response & Proof bytes & Prove s & Verify s\\\midrule\endhead
\midrule\multicolumn{5}{r}{Continued on next page}\\\endfoot
\bottomrule\endlastfoot
train1 & Aggregate & 4,310 & 2.980 & 3.477\\
train1 & Individual & 3,475,990 & 3.048 & 472.659\\
train2 & Aggregate & 4,310 & 0.459 & 1.074\\
train2 & Individual & 539,750 & 0.507 & 73.646\\
\end{longtable}
\endgroup


The corresponding SDA/Per-field public controls select 805/883 pages on train1 and 125/143 pages on train2. Their local request times are 6.437/7.023 s and 1.489/1.683 s. Last-byte ciphertext mutation is rejected in both segment controls.

\section{C3 Complete Recovery Measurements}

\begingroup\small\setlength{\tabcolsep}{5pt}
\begin{longtable}{@{}lrlllrrr@{}}
\caption{Complete-Record Success and Pointwise 95 Percent Wilson Intervals}\label{tab:key-c3-statistics}\\
\toprule Segment & Loss \% & Policy & Protection & Successes & Rate \% & Lower \% & Upper \%\\\midrule\endfirsthead
\toprule Segment & Loss \% & Policy & Protection & Successes & Rate \% & Lower \% & Upper \%\\\midrule\endhead
\midrule\multicolumn{8}{r}{Continued on next page}\\\endfoot
\bottomrule\endlastfoot
train1 & 0.0 & Exact & Dictionary & 30/30 & 100.0 & 88.65 & 100.00\\
train1 & 0.0 & Exact & Uniform & 30/30 & 100.0 & 88.65 & 100.00\\
train1 & 0.0 & Exact & FrequencyD & 30/30 & 100.0 & 88.65 & 100.00\\
train1 & 0.0 & Exact & NoEC & 30/30 & 100.0 & 88.65 & 100.00\\
train1 & 0.0 & Conservative & Dictionary & 30/30 & 100.0 & 88.65 & 100.00\\
train1 & 0.0 & Conservative & Uniform & 30/30 & 100.0 & 88.65 & 100.00\\
train1 & 0.0 & Conservative & FrequencyD & 30/30 & 100.0 & 88.65 & 100.00\\
train1 & 0.0 & Conservative & NoEC & 30/30 & 100.0 & 88.65 & 100.00\\
train1 & 0.1 & Exact & Dictionary & 238/250 & 95.2 & 91.80 & 97.23\\
train1 & 0.1 & Exact & Uniform & 235/250 & 94.0 & 90.34 & 96.33\\
train1 & 0.1 & Exact & FrequencyD & 238/250 & 95.2 & 91.80 & 97.23\\
train1 & 0.1 & Exact & NoEC & 224/250 & 89.6 & 85.20 & 92.80\\
train1 & 0.1 & Conservative & Dictionary & 238/250 & 95.2 & 91.80 & 97.23\\
train1 & 0.1 & Conservative & Uniform & 201/250 & 80.4 & 75.03 & 84.84\\
train1 & 0.1 & Conservative & FrequencyD & 224/250 & 89.6 & 85.20 & 92.80\\
train1 & 0.1 & Conservative & NoEC & 174/250 & 69.6 & 63.64 & 74.97\\
train1 & 1.0 & Exact & Dictionary & 185/250 & 74.0 & 68.23 & 79.04\\
train1 & 1.0 & Exact & Uniform & 153/250 & 61.2 & 55.03 & 67.03\\
train1 & 1.0 & Exact & FrequencyD & 191/250 & 76.4 & 70.76 & 81.24\\
train1 & 1.0 & Exact & NoEC & 114/250 & 45.6 & 39.54 & 51.79\\
train1 & 1.0 & Conservative & Dictionary & 185/250 & 74.0 & 68.23 & 79.04\\
train1 & 1.0 & Conservative & Uniform & 50/250 & 20.0 & 15.51 & 25.40\\
train1 & 1.0 & Conservative & FrequencyD & 97/250 & 38.8 & 32.97 & 44.97\\
train1 & 1.0 & Conservative & NoEC & 17/250 & 6.8 & 4.29 & 10.62\\
train1 & 5.0 & Exact & Dictionary & 37/250 & 14.8 & 10.93 & 19.73\\
train1 & 5.0 & Exact & Uniform & 17/250 & 6.8 & 4.29 & 10.62\\
train1 & 5.0 & Exact & FrequencyD & 43/250 & 17.2 & 13.03 & 22.37\\
train1 & 5.0 & Exact & NoEC & 6/250 & 2.4 & 1.10 & 5.14\\
train1 & 5.0 & Conservative & Dictionary & 27/250 & 10.8 & 7.53 & 15.26\\
train1 & 5.0 & Conservative & Uniform & 0/250 & 0.0 & 0.00 & 1.51\\
train1 & 5.0 & Conservative & FrequencyD & 4/250 & 1.6 & 0.62 & 4.04\\
train1 & 5.0 & Conservative & NoEC & 0/250 & 0.0 & 0.00 & 1.51\\
train2 & 0.0 & Exact & Dictionary & 30/30 & 100.0 & 88.65 & 100.00\\
train2 & 0.0 & Exact & Uniform & 30/30 & 100.0 & 88.65 & 100.00\\
train2 & 0.0 & Exact & FrequencyD & 30/30 & 100.0 & 88.65 & 100.00\\
train2 & 0.0 & Exact & NoEC & 30/30 & 100.0 & 88.65 & 100.00\\
train2 & 0.0 & Conservative & Dictionary & 30/30 & 100.0 & 88.65 & 100.00\\
train2 & 0.0 & Conservative & Uniform & 30/30 & 100.0 & 88.65 & 100.00\\
train2 & 0.0 & Conservative & FrequencyD & 30/30 & 100.0 & 88.65 & 100.00\\
train2 & 0.0 & Conservative & NoEC & 30/30 & 100.0 & 88.65 & 100.00\\
train2 & 0.1 & Exact & Dictionary & 244/250 & 97.6 & 94.86 & 98.90\\
train2 & 0.1 & Exact & Uniform & 239/250 & 95.6 & 92.29 & 97.53\\
train2 & 0.1 & Exact & FrequencyD & 243/250 & 97.2 & 94.33 & 98.64\\
train2 & 0.1 & Exact & NoEC & 232/250 & 92.8 & 88.91 & 95.40\\
train2 & 0.1 & Conservative & Dictionary & 244/250 & 97.6 & 94.86 & 98.90\\
train2 & 0.1 & Conservative & Uniform & 216/250 & 86.4 & 81.60 & 90.10\\
train2 & 0.1 & Conservative & FrequencyD & 233/250 & 93.2 & 89.38 & 95.71\\
train2 & 0.1 & Conservative & NoEC & 192/250 & 76.8 & 71.19 & 81.60\\
train2 & 1.0 & Exact & Dictionary & 185/250 & 74.0 & 68.23 & 79.04\\
train2 & 1.0 & Exact & Uniform & 164/250 & 65.6 & 59.52 & 71.21\\
train2 & 1.0 & Exact & FrequencyD & 185/250 & 74.0 & 68.23 & 79.04\\
train2 & 1.0 & Exact & NoEC & 117/250 & 46.8 & 40.71 & 52.99\\
train2 & 1.0 & Conservative & Dictionary & 185/250 & 74.0 & 68.23 & 79.04\\
train2 & 1.0 & Conservative & Uniform & 65/250 & 26.0 & 20.96 & 31.77\\
train2 & 1.0 & Conservative & FrequencyD & 134/250 & 53.6 & 47.41 & 59.68\\
train2 & 1.0 & Conservative & NoEC & 32/250 & 12.8 & 9.21 & 17.51\\
train2 & 5.0 & Exact & Dictionary & 54/250 & 21.6 & 16.95 & 27.11\\
train2 & 5.0 & Exact & Uniform & 29/250 & 11.6 & 8.20 & 16.16\\
train2 & 5.0 & Exact & FrequencyD & 65/250 & 26.0 & 20.96 & 31.77\\
train2 & 5.0 & Exact & NoEC & 8/250 & 3.2 & 1.63 & 6.19\\
train2 & 5.0 & Conservative & Dictionary & 45/250 & 18.0 & 13.73 & 23.24\\
train2 & 5.0 & Conservative & Uniform & 2/250 & 0.8 & 0.22 & 2.87\\
train2 & 5.0 & Conservative & FrequencyD & 34/250 & 13.6 & 9.90 & 18.40\\
train2 & 5.0 & Conservative & NoEC & 0/250 & 0.0 & 0.00 & 1.51\\
\end{longtable}
\endgroup


\begingroup\small\setlength{\tabcolsep}{5pt}
\begin{longtable}{@{}lrllrrrr@{}}
\caption{Successful and Failed Recovery Time (Milliseconds; Outcome-Specific Medians)}\label{tab:key-c3-time}\\
\toprule Segment & Loss \% & Policy & Protection & $n_s$ & Success ms & $n_f$ & Failure ms\\\midrule\endfirsthead
\toprule Segment & Loss \% & Policy & Protection & $n_s$ & Success ms & $n_f$ & Failure ms\\\midrule\endhead
\midrule\multicolumn{8}{r}{Continued on next page}\\\endfoot
\bottomrule\endlastfoot
train1 & 0.0 & Exact & Dictionary & 30 & 3.693 & 0 & --\\
train1 & 0.0 & Exact & Uniform & 30 & 3.588 & 0 & --\\
train1 & 0.0 & Exact & FrequencyD & 30 & 3.562 & 0 & --\\
train1 & 0.0 & Exact & NoEC & 30 & 3.528 & 0 & --\\
train1 & 0.0 & Conservative & Dictionary & 30 & 8.991 & 0 & --\\
train1 & 0.0 & Conservative & Uniform & 30 & 9.096 & 0 & --\\
train1 & 0.0 & Conservative & FrequencyD & 30 & 8.995 & 0 & --\\
train1 & 0.0 & Conservative & NoEC & 30 & 8.997 & 0 & --\\
train1 & 0.1 & Exact & Dictionary & 238 & 3.866 & 12 & 0.587\\
train1 & 0.1 & Exact & Uniform & 235 & 3.775 & 15 & 0.742\\
train1 & 0.1 & Exact & FrequencyD & 238 & 3.746 & 12 & 0.686\\
train1 & 0.1 & Exact & NoEC & 224 & 3.674 & 26 & 0.687\\
train1 & 0.1 & Conservative & Dictionary & 238 & 9.350 & 12 & 6.760\\
train1 & 0.1 & Conservative & Uniform & 201 & 9.274 & 49 & 3.109\\
train1 & 0.1 & Conservative & FrequencyD & 224 & 9.384 & 26 & 4.708\\
train1 & 0.1 & Conservative & NoEC & 174 & 9.287 & 76 & 3.030\\
train1 & 1.0 & Exact & Dictionary & 185 & 3.974 & 65 & 0.554\\
train1 & 1.0 & Exact & Uniform & 153 & 3.806 & 97 & 0.964\\
train1 & 1.0 & Exact & FrequencyD & 191 & 3.811 & 59 & 0.584\\
train1 & 1.0 & Exact & NoEC & 114 & 3.648 & 136 & 0.670\\
train1 & 1.0 & Conservative & Dictionary & 185 & 9.944 & 65 & 6.421\\
train1 & 1.0 & Conservative & Uniform & 50 & 9.350 & 200 & 2.238\\
train1 & 1.0 & Conservative & FrequencyD & 97 & 9.373 & 153 & 4.206\\
train1 & 1.0 & Conservative & NoEC & 17 & 6.810 & 233 & 1.505\\
train1 & 5.0 & Exact & Dictionary & 37 & 4.298 & 213 & 0.414\\
train1 & 5.0 & Exact & Uniform & 17 & 4.368 & 233 & 0.516\\
train1 & 5.0 & Exact & FrequencyD & 43 & 4.402 & 207 & 0.581\\
train1 & 5.0 & Exact & NoEC & 6 & 3.602 & 244 & 0.321\\
train1 & 5.0 & Conservative & Dictionary & 27 & 11.933 & 223 & 5.810\\
train1 & 5.0 & Conservative & Uniform & 0 & -- & 250 & 0.700\\
train1 & 5.0 & Conservative & FrequencyD & 4 & 8.021 & 246 & 3.050\\
train1 & 5.0 & Conservative & NoEC & 0 & -- & 250 & 0.417\\
train2 & 0.0 & Exact & Dictionary & 30 & 3.507 & 0 & --\\
train2 & 0.0 & Exact & Uniform & 30 & 3.483 & 0 & --\\
train2 & 0.0 & Exact & FrequencyD & 30 & 3.457 & 0 & --\\
train2 & 0.0 & Exact & NoEC & 30 & 3.400 & 0 & --\\
train2 & 0.0 & Conservative & Dictionary & 30 & 7.205 & 0 & --\\
train2 & 0.0 & Conservative & Uniform & 30 & 7.002 & 0 & --\\
train2 & 0.0 & Conservative & FrequencyD & 30 & 7.055 & 0 & --\\
train2 & 0.0 & Conservative & NoEC & 30 & 6.886 & 0 & --\\
train2 & 0.1 & Exact & Dictionary & 244 & 3.575 & 6 & 0.571\\
train2 & 0.1 & Exact & Uniform & 239 & 3.490 & 11 & 1.163\\
train2 & 0.1 & Exact & FrequencyD & 243 & 3.450 & 7 & 0.499\\
train2 & 0.1 & Exact & NoEC & 232 & 3.425 & 18 & 1.075\\
train2 & 0.1 & Conservative & Dictionary & 244 & 7.007 & 6 & 4.599\\
train2 & 0.1 & Conservative & Uniform & 216 & 6.818 & 34 & 2.683\\
train2 & 0.1 & Conservative & FrequencyD & 233 & 6.755 & 17 & 3.832\\
train2 & 0.1 & Conservative & NoEC & 192 & 6.621 & 58 & 2.133\\
train2 & 1.0 & Exact & Dictionary & 185 & 3.647 & 65 & 0.552\\
train2 & 1.0 & Exact & Uniform & 164 & 3.607 & 86 & 0.726\\
train2 & 1.0 & Exact & FrequencyD & 185 & 3.584 & 65 & 0.629\\
train2 & 1.0 & Exact & NoEC & 117 & 3.441 & 133 & 0.685\\
train2 & 1.0 & Conservative & Dictionary & 185 & 7.654 & 65 & 4.795\\
train2 & 1.0 & Conservative & Uniform & 65 & 5.673 & 185 & 1.695\\
train2 & 1.0 & Conservative & FrequencyD & 134 & 6.450 & 116 & 4.866\\
train2 & 1.0 & Conservative & NoEC & 32 & 5.292 & 218 & 1.149\\
train2 & 5.0 & Exact & Dictionary & 54 & 4.074 & 196 & 0.461\\
train2 & 5.0 & Exact & Uniform & 29 & 3.785 & 221 & 0.598\\
train2 & 5.0 & Exact & FrequencyD & 65 & 4.067 & 185 & 0.538\\
train2 & 5.0 & Exact & NoEC & 8 & 3.414 & 242 & 0.355\\
train2 & 5.0 & Conservative & Dictionary & 45 & 7.079 & 205 & 3.892\\
train2 & 5.0 & Conservative & Uniform & 2 & 6.228 & 248 & 0.613\\
train2 & 5.0 & Conservative & FrequencyD & 34 & 6.885 & 216 & 4.223\\
train2 & 5.0 & Conservative & NoEC & 0 & -- & 250 & 0.395\\
\end{longtable}
\endgroup


A dash indicates no observations of that outcome. Clean results use one layout and 30 cases per control; nonzero-loss results use five layouts and 250 cases. The successful/failure populations are retained separately.

\begingroup\small\setlength{\tabcolsep}{5pt}
\begin{longtable}{@{}lrlrrrrrr@{}}
\caption{Dictionary Recovery Components for Successful Fault Requests}\label{tab:c3-components}\\
\toprule Segment & Loss \% & Policy & $n_s$ & Downloaded B & Returned B & Decrypt ms & Decode ms & Init ms\\\midrule\endfirsthead
\toprule Segment & Loss \% & Policy & $n_s$ & Downloaded B & Returned B & Decrypt ms & Decode ms & Init ms\\\midrule\endhead
\midrule\multicolumn{9}{r}{Continued on next page}\\\endfoot
\bottomrule\endlastfoot
train1 & 0.1 & Exact & 238 & 323,584 & 323,584 & 1.028 & 2.015 & 0.752\\
train1 & 0.1 & Conservative & 238 & 1,392,640 & 1,392,640 & 4.214 & 2.182 & 0.712\\
train1 & 1.0 & Exact & 185 & 327,680 & 323,584 & 1.020 & 2.043 & 0.674\\
train1 & 1.0 & Conservative & 185 & 1,392,640 & 1,392,640 & 4.230 & 2.139 & 0.681\\
train1 & 5.0 & Exact & 37 & 364,544 & 323,584 & 1.013 & 2.139 & 0.639\\
train1 & 5.0 & Conservative & 27 & 1,413,120 & 1,392,640 & 4.250 & 2.261 & 0.684\\
train2 & 0.1 & Exact & 244 & 315,392 & 315,392 & 0.962 & 1.833 & 0.644\\
train2 & 0.1 & Conservative & 244 & 1,040,384 & 1,040,384 & 3.060 & 1.907 & 0.627\\
train2 & 1.0 & Exact & 185 & 319,488 & 315,392 & 0.957 & 1.856 & 0.619\\
train2 & 1.0 & Conservative & 185 & 1,040,384 & 1,040,384 & 3.090 & 1.898 & 0.622\\
train2 & 5.0 & Exact & 54 & 346,112 & 296,960 & 0.958 & 1.898 & 0.604\\
train2 & 5.0 & Conservative & 45 & 704,512 & 671,744 & 2.110 & 1.948 & 0.607\\
\end{longtable}
\endgroup


Downloaded bytes include originals and repair shards fetched by the protected reader. Returned bytes are ciphertext objects supplied to the record-recovery layer. Sidecar initialization precedes the recovery timer. Decrypt time includes the checks performed inside the opening helper. Component medians need not sum to the median total.

\Needspace{30\baselineskip}
\subsection{Paired Recovery Outcomes by Layout}
Table~\ref{tab:paired-layouts} pairs the two reported outcomes for each
segment, layout seed, loss fraction, target, and case. FrequencyD and
Dictionary have gains and losses in both directions on train1 at 5\% loss,
whereas Exact never loses a case to Conservative with the Dictionary sidecar
held fixed. Exact and Conservative also agree case by case at 0.1\% and 1\%.
The protection strategies may change both the selected pages and their
stripe membership. These counts characterize the complete configurations.

\begin{table}[t]
\centering\small
\caption{Discordant Recovery Outcomes by Layout Seed}
\label{tab:paired-layouts}
\begin{tabular}{@{}lllrrrrr@{}}
\toprule
Segment & Loss & Comparison & 41001 & 41002 & 41003 & 41004 & 41005 \\
\midrule
train1 & 0.1\% & F vs. D & 1/3 & 1/0 & 1/0 & 0/0 & 0/0 \\
train1 & 0.1\% & D vs. U & 0/0 & 0/1 & 2/0 & 0/0 & 2/0 \\
train1 & 1\% & F vs. D & 2/0 & 1/3 & 4/0 & 3/2 & 2/1 \\
train1 & 1\% & D vs. U & 7/1 & 9/0 & 7/0 & 6/1 & 7/2 \\
train1 & 5\% & F vs. D & 0/1 & 2/2 & 5/1 & 3/1 & 4/3 \\
train1 & 5\% & D vs. U & 9/0 & 3/1 & 3/3 & 2/0 & 8/1 \\
train1 & 5\% & Exact vs. Cons. & 3/0 & 1/0 & 1/0 & 2/0 & 3/0 \\
\midrule
train2 & 0.1\% & F vs. D & 0/0 & 0/0 & 0/0 & 0/1 & 1/1 \\
train2 & 0.1\% & D vs. U & 1/0 & 2/0 & 2/0 & 0/0 & 1/1 \\
train2 & 1\% & F vs. D & 0/0 & 1/2 & 2/1 & 1/1 & 3/3 \\
train2 & 1\% & D vs. U & 3/0 & 5/2 & 7/3 & 5/2 & 9/1 \\
train2 & 5\% & F vs. D & 2/1 & 4/0 & 7/1 & 2/2 & 2/2 \\
train2 & 5\% & D vs. U & 4/0 & 7/0 & 4/2 & 8/2 & 7/1 \\
train2 & 5\% & Exact vs. Cons. & 1/0 & 2/0 & 2/0 & 2/0 & 2/0 \\
\bottomrule
\end{tabular}

\parbox{0.87\linewidth}{\footnotesize Each cell is the number of cases won
only by the first/second configuration among 50 matched cases. F, D, and U
denote FrequencyD, Dictionary, and Uniform under Exact recovery. Exact vs.
Conservative uses the same Dictionary sidecar for both.}
\end{table}

\section{Protection Ledger}

\begingroup\small\setlength{\tabcolsep}{5pt}
\begin{longtable}{@{}llrrrrr@{}}
\caption{Per-Batch Charged Protection State (Five Layouts; Bytes)}\label{tab:protection}\\
\toprule Batch & Policy & Stripes & Cap & Median & Min & Max\\\midrule\endfirsthead
\toprule Batch & Policy & Stripes & Cap & Median & Min & Max\\\midrule\endhead
\midrule\multicolumn{7}{r}{Continued on next page}\\\endfoot
\bottomrule\endlastfoot
train1 / 2020-07-11 & Dictionary & 88 & 741,126 & 738,545 & 738,534 & 738,560\\
train1 / 2020-07-11 & Uniform & 88 & 741,126 & 738,693 & 738,673 & 738,702\\
train1 / 2020-07-11 & FrequencyD & 88 & 741,126 & 738,596 & 738,590 & 738,616\\
train1 / 2020-07-12 & Dictionary & 88 & 746,543 & 738,557 & 738,542 & 738,562\\
train1 / 2020-07-12 & Uniform & 88 & 746,543 & 738,698 & 738,680 & 738,704\\
train1 / 2020-07-12 & FrequencyD & 88 & 746,543 & 738,608 & 738,594 & 738,611\\
train1 / 2020-07-13 & Dictionary & 45 & 378,645 & 377,640 & 377,638 & 377,640\\
train1 / 2020-07-13 & Uniform & 45 & 378,645 & 377,693 & 377,685 & 377,701\\
train1 / 2020-07-13 & FrequencyD & 45 & 378,645 & 377,634 & 377,632 & 377,635\\
train2 / 2020-07-31 & Dictionary & 6 & 54,645 & 50,561 & 50,560 & 50,563\\
train2 / 2020-07-31 & Uniform & 6 & 54,645 & 50,567 & 50,567 & 50,572\\
train2 / 2020-07-31 & FrequencyD & 6 & 54,645 & 50,549 & 50,549 & 50,549\\
train2 / 2020-08-01 & Dictionary & 71 & 601,284 & 595,880 & 595,878 & 595,891\\
train2 / 2020-08-01 & Uniform & 71 & 601,284 & 595,965 & 595,950 & 595,975\\
train2 / 2020-08-01 & FrequencyD & 71 & 601,284 & 595,873 & 595,872 & 595,885\\
train2 / 2020-08-02 & Dictionary & 85 & 718,564 & 713,358 & 713,350 & 713,372\\
train2 / 2020-08-02 & Uniform & 85 & 718,564 & 713,505 & 713,495 & 713,511\\
train2 / 2020-08-02 & FrequencyD & 85 & 718,564 & 713,394 & 713,392 & 713,408\\
train2 / 2020-08-03 & Dictionary & 48 & 404,577 & 402,802 & 402,802 & 402,804\\
train2 / 2020-08-03 & Uniform & 48 & 404,577 & 402,858 & 402,853 & 402,866\\
train2 / 2020-08-03 & FrequencyD & 48 & 404,577 & 402,801 & 402,799 & 402,801\\
\end{longtable}
\endgroup


All listed charged costs satisfy their batch caps. Each cost includes stored parity, signed sidecar metadata and the 64-byte provisioned key/digest pair. NoEC uses zero sidecar bytes. Base ciphertext, audit tags, base manifest and owner state belong to their separate ledgers.

\section{Evidence and Reproduction}
The LaTeX package contains the summary tables and the plot-source rows under
\path{evidence/current-secure}. A separate experiment archive, identified by
\path{dual_hai_20261008_v3}, records the reduced profile for C1 and the full
profile for C2 and C3. Its source-profile index is
\path{CONTRIBUTION_SOURCES.json}. The archive includes case-level
observations, configuration manifests with input and source hashes, and the
experiment source snapshot. The package README identifies the local archive
and the commands for validating the recorded summaries and starting a rerun
from matching inputs.

The summary verifier checks observation-key uniqueness, sample counts,
summary values, and serialized protection caps. The layout breakdown in
Table~\ref{tab:paired-layouts} uses the archived C3 recovery outcomes.
The two public representation controls reported for C2 are stored in the
archive's mechanism-control observations.

\bibliographystyle{IEEEtran}
\bibliography{IEEEabrv,bstcontrol,references}
\end{document}
